\begin{frame}{Hanoi 규칙}\begin{itemize}\item 한 번에 원판 하나만 이동\item 작은 원판 위에 큰 원판을 놓을 수 없음\item 시작 peg의 $n$개를 목적 peg로 이동\end{itemize}\centering\begin{tikzpicture}[scale=1.25,transform shape]\foreach \x/\n in {0/L,3/M,6/R}{\draw[peg] (\x,0)--(\x,2);\draw[peg] (\x-1,0)--(\x+1,0);\node at (\x,-.35){\n};}\foreach \y/\w in {.15/2,.42/1.5,.69/1}{\draw[fill=AlgoOrange!30] (-\w/2,\y) rectangle (\w/2,\y+.2);}\end{tikzpicture}
\end{frame}
\begin{frame}{작은 $n$부터 관찰}$n=1$은 한 번 이동. $n$개라면 가장 큰 원판을 움직이기 전에 위의 $n-1$개를 임시 peg로 치워야 한다.
\end{frame}
\begin{frame}{세 단계 구조}\begin{enumerate}\item $n-1$: start → temp\item 원판 $n$: start → dest\item $n-1$: temp → dest\end{enumerate}\mission{$H(n,s,d,t)$는 $n$개를 $s$에서 $d$로 옮긴다.}
\end{frame}
\begin{frame}{$n=3$의 핵심 상태}\centering
\only<1|handout:0>{\textbf{초기 상태}: $H(3,L,R,M)$ 시작}
\only<2|handout:0>{\textbf{첫 재귀 완료}: $H(2,L,M,R)$ — 3회 이동}
\only<3|handout:1>{\textbf{전체 완료}: 큰 원판 이동 후 $H(2,M,R,L)$ — 총 7회}
\vfill
\begin{tikzpicture}[scale=1.28,transform shape]
\foreach \x/\n in {0/L,3/M,6/R}{\draw[peg] (\x,0)--(\x,1.8);\draw[peg] (\x-1,0)--(\x+1,0);\node at (\x,-.3){\n};}
\only<1|handout:0>{\draw[disk](-.9,.08)rectangle(.9,.3);\draw[disk](-.65,.32)rectangle(.65,.54);\draw[disk](-.4,.56)rectangle(.4,.78);}
\only<2|handout:0>{\draw[disk](-.9,.08)rectangle(.9,.3);\draw[disk](2.35,.08)rectangle(3.65,.3);\draw[disk](2.6,.32)rectangle(3.4,.54);}
\only<3|handout:1>{\draw[disk](5.1,.08)rectangle(6.9,.3);\draw[disk](5.35,.32)rectangle(6.65,.54);\draw[disk](5.6,.56)rectangle(6.4,.78);}
\end{tikzpicture}
\end{frame}
\begin{frame}{Hanoi: 재귀 호출과 현재 작업}\centering
\begin{tikzpicture}[
  level distance=16mm,
  level 1/.style={sibling distance=76mm},
  level 2/.style={sibling distance=31mm},
  every node/.style={tree node,minimum width=18mm,minimum height=8mm,font=\normalsize},
  edge from parent/.style={draw=AlgoGray,thick}
]
\node{$H(3)$}
 child{node{$H(2)$} child{node{$H(1)$}} child{node{$H(1)$}}}
 child{node{$H(2)$} child{node{$H(1)$}} child{node{$H(1)$}}};
\node[draw=AlgoOrange,rounded corners,fill=AlgoOrange!15,text=AlgoOrangeText,
      font=\small\bfseries,inner sep=4pt] at (0,-1.6) {move disk 3};
\end{tikzpicture}
\vfill 두 $H(2)$ 호출 사이에 현재 작업 \textcolor{AlgoOrangeText}{move disk 3}이 있다.
\par\small 각 $H(2)$ 내부에도 두 재귀 호출 사이의 local move가 있으며, 모든 호출은 서로 다른 실행이다.
\end{frame}
\begin{frame}{Hanoi 점화식과 해}\[T(1)=1,\qquad T(n)=2T(n-1)+1.\]\[T(n)=2^n-1=\Th(2^n).\]$n=30$이면 약 10억 번 이동이다.
\begin{takeaway}
모든 이동 순서를 실제로 생성하거나 출력한다면 출력 크기가 $2^n-1$이므로
$\Omega(2^n)$ 시간이 필요하다.
\end{takeaway}
\muted{최소 이동 횟수 $2^n-1$만 계산하는 문제는 별도다.}
\end{frame}
\begin{frame}{Hanoi Takeaway}재귀 코드는 “작은 문제 2회+현재 이동 1회”를 그대로 표현한다. 정확성은 귀납법과 잘 맞지만, 자연스러운 표현이 효율적이라는 뜻은 아니다.
\vfill\muted{지금까지 호출 구조가 고정되어 있었다. Maze에서는 선택 결과에 따라 탐색 경로가 달라진다.}
\end{frame}
